\documentclass[11pt,a4paper]{article}
% ---------------------------------------------------------------------------
% Preambule pro přípravu na bakalářské státnice (MFF UK, Informatika / UI / SU)
% Sazba: pdflatex (UTF-8, čeština). Build: latexmk -pdf src/main.tex
% ---------------------------------------------------------------------------
\usepackage[T1]{fontenc}
\usepackage[utf8]{inputenc}
\usepackage{lmodern}
\usepackage[czech]{babel}
\usepackage[a4paper,margin=2.3cm]{geometry}

\usepackage{amsmath,amssymb,mathtools}
\usepackage{textcomp}
\usepackage{microtype}
\usepackage{xcolor}
\usepackage[export]{adjustbox}   % klíče max width / max height u \includegraphics
\usepackage{graphicx}
\graphicspath{{../}}             % obrázky z poznámek (relativně ke kořeni repa)
\usepackage{enumitem}
\usepackage{titlesec}
\usepackage[most]{tcolorbox}
\usepackage{booktabs}
\usepackage{multicol}
\usepackage{parskip}
\usepackage{hyperref}

% --- chybějící / nestandardní makra z poznámek (psány pro webový renderer) ---
\newcommand{\gt}{>}
\newcommand{\lt}{<}
\providecommand{\set}[1]{\{\,#1\,\}}
\renewcommand{\set}[1]{\{\,#1\,\}}
\providecommand{\braket}[1]{\left\langle #1 \right\rangle}
\renewcommand{\braket}[1]{\left\langle #1 \right\rangle}
\providecommand{\quotedblbase}{\glqq}
\providecommand{\quotesinglbase}{\glq}

% --- hlubší zanoření seznamů (poznámky jdou až do ~7 úrovní) ---
\setlistdepth{9}
\setlist{leftmargin=1.25em, itemsep=1.2pt, topsep=2.5pt, parsep=0pt}
\setlist[itemize,1]{label=\textbullet}
\setlist[itemize,2]{label=\normalfont\bfseries\textendash}
\setlist[itemize,3]{label=\textperiodcentered}
\setlist[itemize,4]{label=$\ast$}
\setlist[itemize,5]{label=$\cdot$}
\setlist[itemize,6]{label=$\circ$}
\setlist[itemize,7]{label=$\diamond$}
\renewlist{itemize}{itemize}{9}
\setlist[itemize,1]{label=\textbullet}
\setlist[itemize,2]{label=\normalfont\bfseries\textendash}
\setlist[itemize,3]{label=\textperiodcentered}
\setlist[itemize,4]{label=$\ast$}
\setlist[itemize,5]{label=$\cdot$}
\setlist[itemize,6]{label=$\circ$}
\setlist[itemize,7]{label=$\diamond$}
\setlist[itemize,8]{label=$\triangleright$}
\setlist[itemize,9]{label=$-$}

% --- barvy a styl nadpisů ---
\definecolor{mffblue}{RGB}{27,54,93}
\definecolor{mffaccent}{RGB}{0,90,156}
\definecolor{practiceyellow}{RGB}{180,120,0}
\definecolor{gapred}{RGB}{150,30,30}

\titleformat{\section}{\Large\bfseries\color{mffblue}}{\thesection}{0.6em}{}
\titleformat{\subsection}{\large\bfseries\color{mffaccent}}{\thesubsection}{0.6em}{}
\titleformat{\subsubsection}{\normalsize\bfseries}{\thesubsubsection}{0.6em}{}
\titlespacing*{\section}{0pt}{1.4em}{0.6em}

\hypersetup{
  colorlinks=true, linkcolor=mffaccent, urlcolor=mffaccent, citecolor=mffaccent,
  pdftitle={Příprava na bakalářské státnice -- Informatika / UI / Strojové učení},
  pdfauthor={Viktor Číhal}, bookmarksnumbered=true,
}

% --- callout boxy ---
% Procvičit: místo, kde poznámky odkazují na vlastní procvičení (není to mezera).
\newtcolorbox{practicebox}{
  enhanced, breakable, colback=practiceyellow!7, colframe=practiceyellow,
  boxrule=0.7pt, left=6pt, right=6pt, top=4pt, bottom=4pt, arc=2pt,
  before skip=5pt, after skip=5pt,
  overlay first={\node[anchor=north west,font=\bfseries\footnotesize\color{practiceyellow}]
    at ([xshift=6pt,yshift=-2pt]frame.north west){};},
  fonttitle=\bfseries\footnotesize, title={\textsc{Procvičit (poznámky odkazují na vlastní trénink)}}
}
% Mezera / doplněno: téma není v poznámkách (nebo jen jako alternativa).
\newtcolorbox{gapbox}{
  enhanced, breakable, colback=gapred!6, colframe=gapred,
  boxrule=0.8pt, left=6pt, right=6pt, top=4pt, bottom=4pt, arc=2pt,
  before skip=6pt, after skip=6pt,
  fonttitle=\bfseries\footnotesize, title={\textsc{Mezera v poznámkách -- ověř / doplň}}
}
% Info / shrnutí
\newtcolorbox{infobox}[1][]{
  enhanced, breakable, colback=mffaccent!5, colframe=mffaccent,
  boxrule=0.7pt, left=6pt, right=6pt, top=4pt, bottom=4pt, arc=2pt,
  before skip=6pt, after skip=6pt, fonttitle=\bfseries, #1
}

\setcounter{tocdepth}{2}
\setcounter{secnumdepth}{3}
\renewcommand{\arraystretch}{1.25}

% --- méně přetékajících řádků u dlouhých kódových úseků ---
\emergencystretch=3em
\makeatletter
\g@addto@macro\ttfamily{\hyphenchar\font=`\-}  % dovol dělení i ve strojopisu
\makeatother
\hbadness=10000  % netiskni varování o podtečení


% ---- doplňkové balíčky pro obrázky (pilot) ----
\usepackage{tikz}
\usetikzlibrary{positioning,arrows.meta,calc,backgrounds}
\usepackage{pgfplots}
\pgfplotsset{compat=1.17}

% ===================== DESIGN: TÉMA vs UČIVO =====================
\definecolor{temagray}{RGB}{90,90,95}
\definecolor{defblue}{RGB}{0,90,156}
\definecolor{intugreen}{RGB}{40,120,70}

% TÉMA = doslovné znění oficiálního požadavku (anchor / ground truth)
\newtcolorbox{tema}{enhanced, breakable, sharp corners=south,
  colback=temagray!8, colframe=temagray, boxrule=0pt, leftrule=3pt,
  left=8pt, right=8pt, top=4pt, bottom=5pt, before skip=12pt, after skip=2pt,
  fonttitle=\bfseries\footnotesize\color{temagray},
  title={\textsc{Téma}\, \textnormal{\footnotesize(oficiální požadavek -- co musíš umět)}}}

% UČIVO = naše vysvětlení (přepsané pro pochopení, ne z poznámek doslova)
\newcommand{\ucivo}{\par\vspace{1pt}\noindent{\footnotesize\itshape\color{intugreen!60!black}Učivo (jak na to):}\par\nopagebreak\vspace{1pt}}

\newtcolorbox{defbox}[1]{enhanced, breakable, colback=defblue!6, colframe=defblue,
  boxrule=0.6pt, arc=1.5pt, left=6pt, right=6pt, top=2pt, bottom=2pt,
  before skip=4pt, after skip=4pt,
  fonttitle=\bfseries\footnotesize, coltitle=defblue, title={Definice: #1}}

\newtcolorbox{intuice}{enhanced, breakable, colback=intugreen!7, colframe=intugreen,
  boxrule=0.6pt, arc=1.5pt, left=6pt, right=6pt, top=2pt, bottom=2pt,
  before skip=4pt, after skip=4pt,
  fonttitle=\bfseries\footnotesize, coltitle=intugreen, title={Intuice}}

\newcommand{\klic}[1]{\par\smallskip\noindent\textbf{\color{defblue}$\blacktriangleright$\ Klíč:}~#1\par\smallskip}

% odpovědní žebřík: co stačí na 1 / 2 / 3 body
\newcommand{\bod}[1]{\tikz[baseline=-0.55ex]\node[circle,fill=defblue,text=white,inner sep=1.3pt,font=\bfseries\scriptsize]{#1};}
\newtcolorbox{odpoved}{enhanced, breakable, colback=defblue!4, colframe=defblue!60,
  boxrule=0.6pt, arc=1.5pt, left=6pt, right=6pt, top=2pt, bottom=3pt,
  before skip=4pt, after skip=4pt, fonttitle=\bfseries\footnotesize, coltitle=defblue,
  title={Odpověď podle bodů \textnormal{\footnotesize(cíl pro průchod: zelená = 2 body)}}}
% aktivní vybavování
\newtcolorbox{recall}{enhanced, breakable, colback=orange!7, colframe=orange!75!black,
  boxrule=0.6pt, arc=1.5pt, left=6pt, right=6pt, top=2pt, bottom=2pt,
  before skip=4pt, after skip=4pt, fonttitle=\bfseries\footnotesize, coltitle=orange!60!black,
  title={Vyzkoušej se \textnormal{\footnotesize(zakryj text a odpověz nahlas)}}}
\newcommand{\past}[1]{\par\smallskip\noindent\textbf{\color{gapred}Pozor (častá past):}~#1\par\smallskip}
% stavový štítek tématu (cíl zvládnutí)
\newcommand{\cil}[1]{\hfill{\footnotesize\colorbox{intugreen!18}{\textbf{cíl: #1}}}}

\begin{document}

{\large\bfseries\color{mffblue} Ukázka nového formátu (pilot)}\par
\vspace{2pt}
{\small Dvě témata na ukázku: jak oddělit \textsc{Téma} (oficiální požadavek) od \emph{Učiva} (naše stručné, intuitivní vysvětlení), s obrázkem tam, kde pomůže.}\par
\vspace{6pt}\hrule\vspace{10pt}

% ============================================================
\subsection*{Část IV / Strojové učení -- Metoda podpůrných vektorů (SVM)}
\noindent{\footnotesize\colorbox{intugreen!18}{\textbf{Priorita: vysoká}} \colorbox{defblue!12}{\textbf{cíl zvládnutí: zelená (2 body)}} \quad specializace = klíčová pro průchod (7/12)}\par\smallskip

\begin{tema}
\textbf{Metoda podpůrných vektorů:} klasifikátor pro lineárně separabilní třídy;
klasifikátor pro lineárně neseparabilní třídy; jádrové funkce; klasifikace do více tříd.
\end{tema}

\ucivo
\begin{intuice}
Mezi dvě třídy chceme proložit \textbf{co nejširší ulici} (pruh bez bodů). Hranice
vede středem ulice. Čím širší ulice, tím lépe klasifikátor generalizuje na nová data.
\end{intuice}

\begin{center}
\begin{tikzpicture}[scale=0.92]
  \begin{scope}
    % ulice (margin)
    \fill[defblue!7] (-0.2,4.2) -- (4.2,-0.2) -- (4.2,1.8) -- (1.8,4.2) -- cycle;
  \end{scope}
  % dělicí nadrovina x+y=4 a okraje x+y=3, x+y=5
  \draw[very thick] (0,4) -- (4,0) node[below right,font=\footnotesize]{$w^\top x + b = 0$};
  \draw[dashed] (-0.2,3.2) -- (3.2,-0.2);
  \draw[dashed] (0.8,4.2) -- (4.2,0.8);
  % záporná třída (kolečka) pod ulicí
  \foreach \p in {(0.6,0.6),(1.2,0.5),(0.5,1.2),(1.0,1.0)} \fill[intugreen] \p circle (2.4pt);
  \node[intugreen,font=\footnotesize] at (0.3,2.0) {třída $-1$};
  % kladná třída (čtverečky) nad ulicí
  \foreach \p in {(3.3,3.3),(2.9,3.6),(3.6,3.0),(3.2,2.9)} \fill[defblue] \p +(-2.4pt,-2.4pt) rectangle +(2.4pt,2.4pt);
  \node[defblue,font=\footnotesize] at (3.9,2.2) {třída $+1$};
  % support vectory na okrajích
  \draw[red,thick] (1.5,1.5) circle (4pt); \fill[intugreen] (1.5,1.5) circle (2.4pt);
  \draw[red,thick] (2.5,2.5) +(-2.4pt,-2.4pt) rectangle +(2.4pt,2.4pt); \fill[defblue] (2.5,2.5) +(-2.4pt,-2.4pt) rectangle +(2.4pt,2.4pt);
  \node[red,font=\scriptsize,align=center] at (3.6,4.0) {podpůrné\\ vektory};
  \draw[red,->] (3.2,3.85) to[bend right=10] (2.65,2.6);
  \draw[red,->] (3.0,3.7) to[bend right=20] (1.65,1.62);
  % šířka marginu
  \draw[<->,thick] (1.9,1.1) -- (1.1,1.9) node[midway,sloped,above,font=\scriptsize]{margin $=\frac{2}{\|w\|}$};
\end{tikzpicture}
\end{center}

\begin{odpoved}
\setlength{\parskip}{2pt}
\noindent\bod{1}\ \textbf{Minimum:} SVM hledá \emph{nadrovinu} $w^\top x+b=0$, která dvě třídy
odděluje s \emph{největším marginem}; klasifikace je $\operatorname{sign}(w^\top x+b)$.\par
\noindent\bod{2}\ \textbf{+ k tomu:} \emph{podpůrné vektory} jsou body na okraji marginu a jediné
určují hranici. Neseparabilní data: \emph{měkký margin} se slack proměnnými $\xi_i$ a penaltou
$C\sum\xi_i$ ($C$ = kompromis úzký margin vs.\ chyby). \emph{Jádro} $K(x,x')$ nahradí skalární
součin a oddělí data ve vyšší dimenzi bez explicitní transformace (RBF, polynom).\par
\noindent\bod{3}\ \textbf{Bonus (jen když zbude čas):} úloha $\min\tfrac12\|w\|^2$ za podmínek
$y_i(w^\top x_i+b)\ge 1$; více tříd přes \emph{one-vs-rest} nebo \emph{one-vs-one}.
\end{odpoved}

\begin{recall}
\begin{itemize}[leftmargin=1.2em,itemsep=0pt,topsep=1pt]
  \item Co jsou podpůrné vektory a proč jen ony určují hranici?
  \item Co dělá parametr $C$? Co se stane při $C\to\infty$?
  \item Proč jádro umožní oddělit neseparabilní data, aniž počítáme transformaci?
\end{itemize}
\end{recall}
\past{Velké $C$ není „lepší model“ -- úzký margin se snadno přeučí. A jádro nemění SVM, jen skalární součin.}

% ============================================================
\subsection*{Část I / Teorie grafů -- Toky v sítích}

\begin{tema}
\textbf{Toky v sítích:} definice sítě a toku v ní; existence maximálního toku (bez důkazu);
princip hledání maximálního toku v síti s celočíselnými kapacitami (např.\ Ford-Fulkersonův algoritmus).
\end{tema}

\ucivo
\begin{intuice}
Síť = potrubí. Každá trubka má \textbf{kapacitu} (max průtok). Pumpujeme ze zdroje $s$ do
spotřebiče $t$ a ptáme se: kolik nejvíc proteče? Omezuje nás vždy nějaké \textbf{úzké hrdlo}.
\end{intuice}

\begin{center}
\begin{tikzpicture}[>=Stealth, node distance=15mm and 20mm,
   every node/.style={circle,draw,minimum size=7mm,font=\small,inner sep=0pt}]
  \node (s) {$s$};
  \node (a) [above right=of s] {$a$};
  \node (b) [below right=of s] {$b$};
  \node (t) [below right=of a] {$t$};
  \tikzset{every edge/.style={draw,->,font=\scriptsize}}
  \draw (s) edge node[above left]{$3$} (a);
  \draw (s) edge node[below left]{$2$} (b);
  \draw (a) edge node[left]{$1$} (b);
  \draw (a) edge node[above right]{$2$} (t);
  \draw (b) edge node[below right]{$3$} (t);
\end{tikzpicture}\\[-2pt]
{\scriptsize Čísla = kapacity hran. Maximální tok zde je 4 (úzké hrdlo: hrany do $t$ nebo ze $s$).}
\end{center}

\begin{defbox}{síť a tok}
\textbf{Síť}: orientovaný graf se zdrojem $s$, spotřebičem $t$ a kapacitami $c(e)\ge 0$.
\textbf{Tok} $f(e)$ splňuje: $0\le f(e)\le c(e)$ (kapacita) a v každém vrcholu mimo $s,t$ se
\emph{kolik přiteče = kolik odteče} (zachování). Velikost toku $|f|$ = čistý odtok ze $s$.
\end{defbox}

\begin{itemize}[leftmargin=1.2em]
  \item \textbf{Existence:} maximální tok vždy existuje (velikost je shora omezená).
  \item \textbf{Ford-Fulkerson (princip):} dokud v \emph{rezervní (reziduální) síti} existuje
  cesta ze $s$ do $t$ (zlepšující cesta), pošli po ní co nejvíc toku; opakuj. U
  \textbf{celočíselných} kapacit algoritmus skončí a dá celočíselný maximální tok.
  \item \textbf{Min-cut = max-flow:} velikost maximálního toku se rovná kapacitě nejmenšího
  řezu (oddělujícího $s$ od $t$). To je to úzké hrdlo.
\end{itemize}
\klic{tok = kapacita + zachování; Ford-Fulkerson = opakuj zlepšující cesty v reziduální síti; max tok = min řez.}

\end{document}
